\documentclass[11pt,a4paper]{article}

% ── packages ──────────────────────────────────────────────────────────────────
\usepackage[margin=2.5cm]{geometry}
\usepackage{amsmath,amssymb}
\usepackage{booktabs}
\usepackage{graphicx}
\usepackage{hyperref}
\usepackage{natbib}
\usepackage{microtype}
\usepackage{parskip}
\usepackage{caption}
\usepackage{float}
\usepackage{xcolor}
\usepackage{array}
\usepackage{multirow}

\hypersetup{
  colorlinks=true,
  linkcolor=black,
  citecolor=black,
  urlcolor=blue
}

\captionsetup{font=small, labelfont=bf, width=0.95\textwidth}

% ── title ─────────────────────────────────────────────────────────────────────
\title{\textbf{Bilateral Symmetry as a Filter-Robust Discriminator\\
Between Real and GAN-Generated Face Images}}

\author{Scott Cundill\\
\small Independent Researcher, Okinawa, Japan\\
\small \href{mailto:scott.cam@gmail.com}{scott.cam@gmail.com}}

\date{}

% ══════════════════════════════════════════════════════════════════════════════
\begin{document}
\maketitle

% ── abstract ──────────────────────────────────────────────────────────────────
\begin{abstract}
Generative adversarial networks now produce synthetic face images that are,
in many cases, indistinguishable from real photographs by human observers.
Detection methods fall into two broad camps: deep-learning classifiers, and
handcrafted-feature approaches that target specific, interpretable geometric
artifacts. We report a measurement of global bilateral symmetry on 1,000 real
face images from the Flickr-Faces-HQ dataset and 1,000 GAN-generated images
from a StyleGAN release, under three independent image-preprocessing filters:
raw greyscale, topographic isoline, and edge.
GAN-generated faces show significantly higher bilateral symmetry than real
faces under all three filters, with effect sizes of $d = 0.20$--$0.28$ and
$p < 10^{-4}$ throughout. The result is independently confirmed on a held-out
validation sample of 1,000 images drawn from a separate partition of the same
dataset ($d = 0.14$--$0.40$, $p < 0.002$ across all three filters). A parallel
analysis on 1,000 diffusion-generated faces (DALL-E~3, FLUX1, SDXL) reveals a
qualitatively different pattern: diffusion faces show no global luminance
symmetry bias but are dramatically more symmetric in their edge structure
($d = 1.06$, $p = 2.5 \times 10^{-108}$). The result persists under an edge
filter that suppresses global lighting and pose information, substantially
weakening the principal confounds for whole-image symmetry measurement. Real
faces additionally show greater multifractal width, Shannon entropy, and
junction complexity. Global bilateral symmetry is a simple, parameter-free,
interpretable feature whose filter-resolved behaviour reveals architecturally
distinct geometric signatures in GAN and diffusion generators, suggesting
utility as a component in ensemble detectors for synthetic face images.
\end{abstract}

% ── 1. Introduction ───────────────────────────────────────────────────────────
\section{Introduction}

Generative adversarial networks now produce synthetic face images that are, in
many cases, indistinguishable from photographs by human observers
\citep{goodfellow2014generative, karras2019style}. Detection methods fall into
two broad camps: deep-learning classifiers trained end-to-end on real and
synthetic pairs \citep{tolosana2022survey}, and handcrafted-feature approaches
that target specific, interpretable geometric artifacts. The latter are of
particular interest because, when a feature discriminates, the reason is
legible.

Recent handcrafted work has focused on local facial asymmetries, specifically
differences between the two eyes in colour \citep{matern2019exploiting},
corneal highlight patterns \citep{hu2021exposing}, and iris structure. The
underlying observation is sound: GANs impose no explicit symmetry constraint
during generation, so bilateral consistency emerges only as a statistical
property of the training data, not as an architectural guarantee
\citep{gragnaniello2022detection}.

Whether this extends to the \emph{global} bilateral symmetry of the whole face
has not, to our knowledge, been reported. This is the question we address.
Global bilateral symmetry---the correspondence between the left and right halves
of the full image---is a single scalar measurement requiring no landmark
detection or segmentation. It is computationally trivial, and its behaviour
under different image representations has not been tested in this context.

We measured global bilateral symmetry on 1,000 real faces from FFHQ
\citep{karras2019style} and 1,000 GAN-generated faces from the 140k Real and
Fake Faces dataset \citep{tunguz2020, xhlulu2020}, under three independent
preprocessing filters: raw greyscale, topographic isoline, and edge.
GAN-generated faces are measurably more bilaterally symmetric than real faces
under all three, including the edge filter, which suppresses global pose and
lighting---the principal confounds for any whole-image symmetry measurement.
Three secondary metrics (multifractal width, Shannon entropy, skeleton junction
count) show consistent separation under the raw filter.

We make no claim of universal separability between real and synthetic faces.
The result is specific to these datasets and this metric. Its value is that it
is simple, interpretable, and stable across basic changes in image
representation. This work arose from the development of SCOTT (Structural and
Coherent Order in Topological Transforms), an open-source framework for
measuring geometric and topological order in 2D images, in the course of which
a battery of symmetry and topology metrics was applied to publicly available
face image datasets. The bilateral symmetry finding reported here emerged from
that broader measurement programme.

% ── 2. Datasets ───────────────────────────────────────────────────────────────
\section{Datasets}

Real face images were drawn from the Flickr-Faces-HQ dataset (FFHQ), collected
by NVIDIA \citep{karras2019style}. GAN-generated images were drawn from the ``1
Million Fake Faces'' dataset \citep{tunguz2020}, a StyleGAN release compiled
alongside FFHQ into the ``140k Real and Fake Faces'' package on Kaggle
\citep{xhlulu2020}. From each source, 1,000 images were selected, giving a
balanced sample of $n = 2{,}000$. Both subsets were drawn from the test
partition of the Kaggle dataset, held separate from any training split by the
dataset's compiler. FFHQ images are independent photographs of distinct
individuals; StyleGAN images are independently sampled from the generator's
latent space and share no identity information. All images were converted to
greyscale prior to analysis.

The Kaggle dataset stores images at $256 \times 256$ pixel resolution. The SCOTT
pipeline normalises all inputs to $512 \times 512$ before analysis using cubic
spline interpolation (\texttt{cv2.INTER\_CUBIC}), to ensure uniform filter radii
and consistent spatial frequency response across all datasets regardless of
original resolution. This represents a $2\times$ upscaling of the original data.
Because the bilateral symmetry metric and the secondary metrics operate on
pixel-value relationships, upscaling affects both populations equally,
preserving the relative discrimination between them. The SFHQ-T2I diffusion
images are natively $1024 \times 1024$ and are downscaled to $512 \times 512$
by the same normalisation step. All datasets were processed consistently under
pipeline version v16.2.

% ── 3. Methods ────────────────────────────────────────────────────────────────
\section{Methods}

\subsection{Bilateral symmetry metric}

Global bilateral symmetry was computed as follows. Each greyscale image was
normalised to the range $[0, 1]$ as a floating-point array, then flipped
horizontally to produce its mirror image. Bilateral symmetry is defined as:

\begin{equation}
  \mathrm{BS} = 1 - \mathrm{MAD}(I, F)
\end{equation}

where $I$ is the normalised image ($512 \times 512$, float $[0,1]$), $F$ is its
horizontal flip, and $\mathrm{MAD}$ is the mean absolute difference:

\begin{equation}
  \mathrm{MAD}(I, F) = \frac{1}{HW} \sum_{h=1}^{H} \sum_{w=1}^{W}
  \bigl| I(h,w) - F(h,w) \bigr|
\end{equation}

The result ranges from 0 (no mirror correspondence) to 1.0 (perfect bilateral
symmetry). This is a direct pixel-wise deviation from mirror symmetry, not a
normalised cross-correlation. The MAD formulation penalises absolute pixel
differences uniformly across the image, making it a stringent and interpretable
measure of global mirror correspondence. A secondary measure, left-right
brightness imbalance, was computed as the absolute difference in mean pixel
intensity between the left and right halves of the image.

\subsection{Image preprocessing filters}

Three independent preprocessing filters were applied to each image before metric
computation. Each suppresses a different class of image information, providing
three partially independent tests of the same hypothesis. A metric that responds
consistently across all three representations is unlikely to reflect a single
artefact type.

\textit{Raw.} The greyscale image at $512 \times 512$ with no further
processing. This is the control representation, retaining all spatial
information including lighting gradients, texture, and fine detail.

\textit{Topographic isoline.} Contrast-Limited Adaptive Histogram Equalisation
(CLAHE, clip limit 4, tile grid $8 \times 8$) was first applied to normalise
local contrast. The image was then smoothed with a bilateral filter (diameter 9,
$\sigma = 75$ in both colour and coordinate space). Intensity isolines were
extracted at uniform intervals of 15 grey levels across the full 0--255 range,
producing 16 contour levels drawn as 1-pixel lines on a black canvas. The result
suppresses texture and retains only the large-scale geometric structure of the
brightness landscape. All topographic-filtered data were processed under
pipeline version v16.2 (ingestion timestamps: 2026-06-01 to 2026-06-03),
ensuring consistent CLAHE parameters across all datasets.

\textit{Edge.} The greyscale image was smoothed with a $5 \times 5$ Gaussian
kernel, then passed through a Canny edge detector with thresholds of 30 (low)
and 120 (high), chosen to retain fine facial detail including shallow surface
gradients. A $3 \times 3$ rectangular closing operation was applied to bridge
small gaps in detected edges. This representation responds to local gradient
magnitude and is largely insensitive to global brightness gradients and
low-frequency composition.

\subsection{Secondary metrics}

Three secondary metrics were computed on the raw filter only, as supporting
observations rather than primary claims.

\textit{Multifractal width} measures the spread of the multifractal spectrum of
the greyscale image, capturing the richness of intensity variation across spatial
scales \citep{chhabra1989}.

\textit{Shannon entropy} of the greyscale histogram measures the distribution of
pixel intensities; higher entropy indicates a more uniform spread and, in natural
images, greater tonal complexity \citep{shannon1948}.

\textit{Junction count} is the number of branch-point junctions in the binary
skeleton of the image, extracted after thresholding. It reflects the topological
complexity of the intensity surface.

\subsection{Statistical procedure}

Group means were compared using Welch's two-sample $t$-test, which does not
assume equal variances. Tests were two-sided. Exact $p$-values are reported
throughout. Effect size is reported as Cohen's $d$, computed as the mean
difference divided by the pooled standard deviation. Each filter condition was
treated as a separate planned comparison. No correction for multiple comparisons
was applied across the three filter tests, as they address the same hypothesis
under different image representations rather than testing independent hypotheses.
A Bonferroni-corrected threshold of $p < 0.017$ would nonetheless be met by all
three results.

For the three secondary metrics, a Benjamini--Hochberg correction at FDR $= 0.05$
was applied; all three remain significant after correction (adjusted $p$-values:
multifractal width $1.2 \times 10^{-6}$, Shannon entropy $2.2 \times 10^{-3}$,
junction count $1.5 \times 10^{-2}$).

% ── 4. Results ────────────────────────────────────────────────────────────────
\section{Results}

\subsection{Bilateral symmetry results}

GAN-generated face images showed consistently higher bilateral symmetry than
real face images across all three filter representations. Table~\ref{tab:main}
reports means, standard deviations, test statistics, effect sizes, and
independent validation results.

\begin{table}[H]
\centering
\caption{Bilateral symmetry (BS) in real versus GAN-generated face images under
three filter representations. $n = 1{,}000$ per group. Val.\ $d$ and Val.\ $p$
report results from an independent validation sample of 1,000 images drawn from
the Kaggle train partition (random seed 42), never used in the main analysis.
Cohen's $d$ is reported as AI minus real throughout.}
\label{tab:main}
\small
\begin{tabular}{lcccccccccc}
\toprule
Filter & $\bar{x}_\text{real}$ & $s_\text{real}$ & $\bar{x}_\text{AI}$ &
$s_\text{AI}$ & $t$ & $p$ & $d$ & Val.\ $d$ & Val.\ $p$ \\
\midrule
Raw         & 0.8008 & 0.0673 & 0.8198 & 0.0700 & $-6.19$ &
  $7.5\times10^{-10}$ & $0.28$ & $0.40$ & $7.9\times10^{-19}$ \\
Topographic & 0.5887 & 0.0289 & 0.5948 & 0.0288 & $-4.79$ &
  $1.8\times10^{-6}$  & $0.21$ & $0.14$ & $1.4\times10^{-3}$  \\
Edge        & 0.8105 & 0.0616 & 0.8225 & 0.0616 & $-4.37$ &
  $1.3\times10^{-5}$  & $0.20$ & $0.29$ & $1.4\times10^{-10}$ \\
\bottomrule
\end{tabular}
\end{table}

All three comparisons are significant at $p < 10^{-4}$. Effect sizes are small
to modest ($d = 0.20$--$0.28$), indicating that the distributions overlap
substantially and the metric alone would not constitute a reliable classifier.
The direction is consistent: GAN-generated faces are more bilaterally symmetric
than real faces under every representation tested. The raw filter produces the
largest effect ($d = 0.28$), as expected, since it retains all image information
including lighting and composition. The topographic and edge filters reduce the
effect size modestly but do not eliminate it.

\subsection{Independent validation}

To confirm the result on unseen data, we drew an independent sample of 1,000
real and 1,000 AI images from the train partition of the Kaggle dataset, which
was not used in the main analysis. Images were selected by random sampling with
seed 42. The bilateral symmetry effect replicates under all three filters, with
$d = 0.40$ ($p = 7.9 \times 10^{-19}$) on the raw filter, $d = 0.29$
($p = 1.4 \times 10^{-10}$) on the edge filter, and $d = 0.14$
($p = 1.4 \times 10^{-3}$) on the topographic filter. Effect sizes are
comparable or larger than in the main analysis, providing no evidence that the
original result was driven by characteristics specific to the valid partition.

\subsection{Edge filter}

The edge filter result warrants specific attention. The Canny edge representation
responds to local gradient magnitude and is largely insensitive to global image
composition, low-frequency lighting gradients, and framing. Pose-related
asymmetry, which affects the low-frequency brightness distribution of the image,
contributes minimally to an edge map dominated by local boundary structure. The
persistence of the bilateral symmetry difference under this representation
($d = 0.20$, $p = 1.3 \times 10^{-5}$) is therefore the central robustness
result of this paper. It does not eliminate the confound entirely, but it
substantially weakens it as an explanation.

\subsection{Secondary metrics}

Table~\ref{tab:secondary} reports three secondary metrics on the raw filter. All
three show consistent separation, with real faces scoring higher than
GAN-generated faces on every measure.

\begin{table}[H]
\centering
\caption{Secondary metrics (raw filter only), real versus GAN-generated faces.
$n = 1{,}000$ per group. Cohen's $d$ is negative throughout, indicating
real $>$ AI. Adjusted $p$-values shown (Benjamini--Hochberg FDR $= 0.05$).}
\label{tab:secondary}
\small
\begin{tabular}{lcccccccc}
\toprule
Metric & $\bar{x}_\text{real}$ & $s_\text{real}$ & $\bar{x}_\text{AI}$ &
$s_\text{AI}$ & $t$ & $p_\text{adj}$ & $d$ & Direction \\
\midrule
Multifractal width & 0.6027 & 0.5103 & 0.4968 & 0.4142 &
  $5.09$ & $1.2\times10^{-6}$ & $-0.23$ & real $>$ AI \\
Shannon entropy    & 7.4206 & 0.3121 & 7.3735 & 0.3473 &
  $3.19$ & $2.2\times10^{-3}$ & $-0.14$ & real $>$ AI \\
Junction count     & 1445.4 & 612.9  & 1377.5 & 637.3  &
  $2.43$ & $1.5\times10^{-2}$ & $-0.11$ & real $>$ AI \\
\bottomrule
\end{tabular}
\end{table}

Real faces show greater multifractal width, higher Shannon entropy, and higher
junction counts than GAN-generated faces, consistent with real photographic
images carrying richer multi-scale complexity than StyleGAN outputs, a pattern
previously observed in frequency-domain analyses \citep{shannon1948,
chhabra1989}. These are reported as supporting observations; the primary claim
of this paper concerns bilateral symmetry. The junction count result in
particular ($d = -0.11$, $p_\text{adj} = 0.015$) is the weakest signal and
should be interpreted cautiously.

\begin{figure}[H]
\centering
\includegraphics[width=\textwidth]{figure1_threeway_violin.png}
\caption{Bilateral symmetry distributions for real (FFHQ), StyleGAN, and
diffusion (SFHQ-T2I: DALL-E~3, FLUX1, SDXL mixed) face images under three
filter representations. Horizontal bars indicate group means.}
\label{fig:violin}
\end{figure}

\begin{figure}[H]
\centering
\includegraphics[width=\textwidth]{figure2_secondary_metrics_bar.png}
\caption{Secondary metrics (raw filter only) for real versus StyleGAN-generated
faces. Error bars show standard error of the mean.}
\label{fig:secondary}
\end{figure}

\subsection{Diffusion-generated faces}

A separate analysis was conducted on 1,000 images drawn at random (seed 42) from
the SFHQ-T2I dataset \citep{beniaguev2024}, a publicly available collection of
diffusion-generated face images spanning three architecturally distinct models:
DALL-E~3, FLUX1, and SDXL. Images are natively $1024 \times 1024$ and were
downscaled to $512 \times 512$ by the standard SCOTT normalisation step. The
dataset is curated and generated from text prompts designed to produce frontal,
well-lit faces; selection biases inherent in the prompting process may affect the
results, and this is acknowledged as a limitation.

\begin{table}[H]
\centering
\caption{Bilateral symmetry: real faces versus diffusion-generated faces
(SFHQ-T2I, $n = 1{,}000$ per group). The edge filter result is the primary
finding. ns = not significant ($p > 0.05$).}
\label{tab:diffusion}
\small
\begin{tabular}{lcccccccc}
\toprule
Filter & $\bar{x}_\text{real}$ & $\bar{x}_\text{diff}$ & $s_\text{real}$ &
$s_\text{diff}$ & $t$ & $p$ & $d$ & Direction \\
\midrule
Raw         & 0.8008 & 0.7962 & 0.0673 & 0.0555 & $1.68$ &
  $0.093$ & $-0.08$ & real $>$ diff (ns) \\
Topographic & 0.5887 & 0.5827 & 0.0289 & 0.0472 & $3.42$ &
  $6.5\times10^{-4}$ & $-0.15$ & real $>$ diff \\
\rowcolor{yellow!20}
Edge        & 0.8105 & 0.8682 & 0.0616 & 0.0465 & $-23.65$ &
  $2.5\times10^{-108}$ & $+1.06$ & diff $>$ real \\
\bottomrule
\end{tabular}
\end{table}

The results reveal a striking dissociation across filter representations. Under
the raw filter, diffusion faces show no significant difference from real faces
($d = -0.08$, $p = 0.093$). Under the topographic filter, they are marginally
less symmetric ($d = -0.15$, $p = 6.5 \times 10^{-4}$). Under the edge filter,
however, diffusion faces are dramatically more symmetric than real faces, with an
effect size of $d = 1.06$, $p = 2.5 \times 10^{-108}$---an effect approximately
five times larger than the StyleGAN result under the same filter.

This dissociation suggests that diffusion models produce a qualitatively
different kind of regularity from StyleGAN. StyleGAN symmetry is primarily a
global, luminance-level phenomenon, visible under raw and topographic
representations. Diffusion symmetry is predominantly local and edge-based: the
boundary geometry of diffusion-generated faces is almost perfectly mirrored in a
way that neither real faces nor StyleGAN faces are. The two generator classes
leave distinct geometric signatures, and the bilateral symmetry metric under
different filter representations separates them.

% ── 5. Discussion ─────────────────────────────────────────────────────────────
\section{Discussion}

\subsection{The symmetry difference}

GAN-generated face images are measurably more bilaterally symmetric than real
photographs, across all three filter representations tested. The effect is small
to modest in magnitude ($d = 0.20$--$0.28$), meaning the distributions overlap
substantially, and the metric alone would not constitute a reliable classifier.
That is not a weakness of the finding---it is an accurate characterisation of it.

The most plausible mechanism is straightforward. StyleGAN was trained on FFHQ, a
dataset of aesthetically curated, roughly frontal, well-lit face photographs.
Curation of this kind selects implicitly for compositional regularity, and the
generator, optimising to reproduce the statistical properties of its training
data, inherits and amplifies that regularity. Real photography, by contrast,
includes asymmetric expression, uneven lighting, slight head rotation, and the
natural bilateral asymmetry of human faces themselves. This explanation is
consistent with the data, but cannot be confirmed by it; the present study does
not manipulate training data composition or image curation.

\subsection{The edge filter as a confound test}

The standard objection to any whole-image symmetry measurement is that the
effect reflects pose or lighting rather than anything intrinsic to the image
source. A face photographed slightly off-centre, or lit from one side, will
score lower on bilateral symmetry regardless of whether it is real or synthetic.
This objection applies most directly to the raw filter result.

The edge filter addresses this directly. The Canny edge representation responds
to local gradient magnitude at pixel boundaries and is largely insensitive to
global brightness gradients and low-frequency composition. Pose-related
asymmetry contributes minimally to an edge map dominated by local boundary
structure. The persistence of the bilateral symmetry difference under the edge
filter ($d = 0.20$, $p = 1.3 \times 10^{-5}$) is therefore the central
robustness result of this paper. It does not eliminate the confound entirely---
strongly off-axis poses would still produce asymmetric edge maps---but it
substantially weakens it as an explanation.

\subsection{Two kinds of regularity}

The diffusion results introduce a distinction the StyleGAN analysis alone could
not reveal. StyleGAN faces are more symmetric than real faces at the global
luminance level. Diffusion faces show no such global regularity, but their edge
structure is dramatically more symmetric, with an effect size five times larger
than the StyleGAN edge result. These are architecturally distinct failure modes.
StyleGAN inherits the compositional regularity of its curated training data.
Diffusion models, generating images through iterative denoising, apparently
impose a strong bilateral constraint on local boundary geometry that does not
manifest at the global luminance level. Whether this reflects a symmetry prior
in the diffusion process, a property of the SFHQ-T2I prompts, or a more general
characteristic of text-to-image face generation is an open question the present
data cannot resolve.

\subsection{Secondary metrics}

Multifractal width, Shannon entropy, and junction count all show real faces
scoring higher than GAN-generated faces under the raw filter. This is consistent
with a broader pattern: real photographic images carry richer multi-scale
complexity than GAN outputs, a difference previously observed in frequency-domain
analyses \citep{shannon1948, chhabra1989}. These are reported as supporting
observations, not independent claims. The junction count result ($d = -0.11$,
$p_\text{adj} = 0.015$) is the weakest signal in the paper and should be
interpreted cautiously.

\subsection{Comparison with existing methods}

Diffusion models have substantially reduced many classic GAN artefacts---spectral
fingerprints, texture inconsistencies, local facial asymmetries---making
traditional handcrafted detectors less reliable on newer generators. The
filter-resolved bilateral symmetry approach described here responds to this
shift: rather than targeting a specific artefact type, it measures a geometric
property whose expression changes across generator classes in a diagnostically
useful way.

Existing handcrafted-feature detection methods focus predominantly on local
asymmetries: per-eye colour differences \citep{matern2019exploiting}, corneal
highlight inconsistencies \citep{hu2021exposing}, and iris structure. These
methods identify specific anatomical locations where GANs fail to impose
consistency. The present metric is complementary rather than competing---it
captures whether the face as a whole is more symmetric than expected, without
localising the source. The two approaches could reasonably be combined in a
multi-feature ensemble, where global bilateral symmetry serves as a fast,
parameter-free first-pass filter. CNN-based classifiers applied to the same 140k
dataset routinely achieve $>95\%$ accuracy \citep{tolosana2022survey}; the
present metric's value is not in outperforming deep classifiers but in being
interpretable, lightweight, and training-free.

\subsection{Limitations}

Several limitations should be noted. First, this study uses specific dataset
pairs: FFHQ for real faces, one StyleGAN release for GAN-generated faces, and
SFHQ-T2I for diffusion-generated faces. Results may not generalise to all
datasets or generator versions within each class---including StyleGAN3/3+ or
other diffusion families (Stable Diffusion, Midjourney). Second, both FFHQ and
the Kaggle sets are heavily curated and predominantly frontal; the effect sizes
reported here may not replicate on datasets with greater pose variation or
in-the-wild conditions. Third, the bilateral symmetry metric is a single scalar
computed over the whole image, is sensitive to strongly off-axis poses, and does
not localise the source of asymmetry. Fourth, the mechanism underlying the
diffusion edge-symmetry result is not established by the present data and
warrants dedicated investigation. Future work should test on more varied pose
distributions, adversarially augmented images, and video-frame inputs.

% ── 6. Conclusion ─────────────────────────────────────────────────────────────
\section{Conclusion}

Global bilateral symmetry, measured as a simple pixel-wise deviation from mirror
correspondence, discriminates GAN-generated face images from real photographs
with small but consistent effect sizes across three independent image
representations, confirmed on a held-out validation sample. The result is robust
under an edge filter that suppresses global lighting and composition
($d = 0.20$, $p = 1.3 \times 10^{-5}$). A parallel analysis on 1,000
diffusion-generated faces (DALL-E~3, FLUX1, SDXL) reveals a qualitatively
different pattern: diffusion faces show no global luminance symmetry bias but
are dramatically more symmetric in their edge structure
($d = 1.06$, $p = 2.5 \times 10^{-108}$). StyleGAN and diffusion models leave
distinct geometric signatures that the bilateral symmetry metric, applied across
filter representations, can separate. These findings suggest that filter-resolved
bilateral symmetry may be a useful component in ensemble detectors for synthetic
face images, with the appropriate filter representation depending on the target
generator class.

% ── 7. Data and code ──────────────────────────────────────────────────────────
\section{Data and code}

Code, measurement data, and analysis scripts are available at:
\url{https://github.com/scottcundill/scott-framework}

The repository contains the SCOTT framework source code, all measurement CSVs
used to produce the tables and figures in this paper, and the reproduction
script. Real face images are available from FFHQ \citep{karras2019style};
GAN-generated images from the 140k Real and Fake Faces dataset
\citep{xhlulu2020}; diffusion images from SFHQ-T2I \citep{beniaguev2024}.

% ── 8. Acknowledgements ───────────────────────────────────────────────────────
\section*{Acknowledgements}

This work arose from the development of SCOTT (Structural and Coherent Order in
Topological Transforms), an open-source geometric image analysis framework, in
the course of which the author applied a battery of topographic and symmetry
metrics to publicly available face image datasets. The bilateral symmetry finding
reported here emerged from that broader measurement programme rather than from a
study designed specifically to test it. The author made extensive use of large
language models---including Claude (Anthropic), GPT-4 (OpenAI), and Gemini
(Google)---for code development, statistical guidance, and manuscript editing.
Interpretation and conclusions are the author's.

% ── References ────────────────────────────────────────────────────────────────
\bibliographystyle{plainnat}

\begin{thebibliography}{99}

\bibitem[Beniaguev(2024)]{beniaguev2024}
Beniaguev, D. (2024).
\newblock Synthetic Faces High Quality -- Text 2 Image (SFHQ-T2I) Dataset.
\newblock Kaggle. \url{https://www.kaggle.com/datasets/selfishgene/sfhq-t2i-synthetic-faces-from-text-2-image-models}
\newblock DOI: 10.34740/kaggle/dsv/9548853

\bibitem[Chhabra \& Jensen(1989)]{chhabra1989}
Chhabra, A., \& Jensen, R.~V. (1989).
\newblock Direct determination of the $f(\alpha)$ singularity spectrum.
\newblock \textit{Physical Review Letters}, 62(12), 1327--1330.

\bibitem[Goodfellow et~al.(2014)]{goodfellow2014generative}
Goodfellow, I., Pouget-Abadie, J., Mirza, M., Xu, B., Warde-Farley, D.,
  Ozair, S., Courville, A., \& Bengio, Y. (2014).
\newblock Generative adversarial nets.
\newblock \textit{Advances in Neural Information Processing Systems}, 27,
  2672--2680.

\bibitem[Gragnaniello et~al.(2022)]{gragnaniello2022detection}
Gragnaniello, D., Marra, F., \& Verdoliva, L. (2022).
\newblock Detection of AI-generated synthetic faces.
\newblock In \textit{Handbook of Digital Face Manipulation and Detection}.
  Springer International Publishing, pp.\ 191--212.

\bibitem[Hu et~al.(2021)]{hu2021exposing}
Hu, S., Li, Y., \& Lyu, S. (2021).
\newblock Exposing GAN-generated faces using inconsistent corneal specular
  highlights.
\newblock \textit{ICASSP 2021}, pp.\ 2500--2504.
\newblock DOI: 10.1109/ICASSP39728.2021.9414582

\bibitem[Karras et~al.(2019)]{karras2019style}
Karras, T., Laine, S., \& Aila, T. (2019).
\newblock A style-based generator architecture for generative adversarial
  networks.
\newblock \textit{Proceedings of the IEEE/CVF CVPR}, pp.\ 4401--4410.

\bibitem[Matern et~al.(2019)]{matern2019exploiting}
Matern, F., Riess, C., \& Stamminger, M. (2019).
\newblock Exploiting visual artifacts to expose deepfakes and face
  manipulations.
\newblock \textit{IEEE WACV Workshops}, pp.\ 83--92.
\newblock DOI: 10.1109/WACVW.2019.00020

\bibitem[Shannon(1948)]{shannon1948}
Shannon, C.~E. (1948).
\newblock A mathematical theory of communication.
\newblock \textit{Bell System Technical Journal}, 27(3), 379--423.

\bibitem[Tolosana et~al.(2022)]{tolosana2022survey}
Tolosana, R., Vera-Rodriguez, R., Fierrez, J., Morales, A., \&
  Ortega-Garcia, J. (2022).
\newblock GAN-generated faces detection: A survey and new perspectives.
\newblock \textit{arXiv preprint arXiv:2202.07145}.

\bibitem[Tunguz(2020)]{tunguz2020}
Tunguz, B. (2020).
\newblock 1 Million Fake Faces.
\newblock Kaggle. \url{https://www.kaggle.com/datasets/tunguz/1-million-fake-faces}

\bibitem[xhlulu(2020)]{xhlulu2020}
xhlulu (2020).
\newblock 140k Real and Fake Faces.
\newblock Kaggle. \url{https://www.kaggle.com/datasets/xhlulu/140k-real-and-fake-faces}

\end{thebibliography}

\end{document}
